\bookchapter{Listening to the GameController}\label{k:gc}

In a game the robots get their orders from the referee through the \textbf{GameController}: a program on the referee's laptop that broadcasts the game state twice per second. A robot that ignores it is a robot that walks during \emph{set}, plays while penalised, and gets taken off the field. Each robot must also send a status packet back. \rtr{The GameController protocol} explains the protocol and shows a parser in Rust; this chapter writes the same in Python, wires it into the blackboard, and tests it against the real GameController. Milestone: one.

\begin{needbox}
\begin{itemize}[nosep]
  \item The loops and the debug tool (Chapters~\ref{k:loops} and~\ref{k:rc}). The official GameController for the HSL, running on the laptop.
  \item Background: \rtr{The GameController protocol}, \rtr{Game flow as a state machine}, \rtr{UDP Networking}.
\end{itemize}
\end{needbox}

\section*{Parsing with \texttt{struct}}
The control message has $158$ bytes: an $18$-byte header and two $70$-byte team blocks; the status message has $32$ bytes (layouts in the concepts book, protocol version $20$ and $4$). Python's \texttt{struct} module turns such layouts into one line each: \texttt{<} means little-endian with no padding, \texttt{B} a byte, \texttt{H}/\texttt{h} an unsigned/signed $16$-bit number, \texttt{f} a $32$-bit float, \texttt{4s} four raw bytes.
\begin{verbatim}
import struct
from dataclasses import dataclass

CONTROL = struct.Struct("<4s10B2h")     # 18-byte header
TEAM = struct.Struct("<6B2H60B")        # 70-byte team block
RETURN = struct.Struct("<4s4B3ff2f")    # 32-byte status message
STATES = ["initial", "ready", "set", "playing", "finished"]

@dataclass
class GameState:
    packet: int
    state: str
    stopped: bool
    set_play: int
    first_half: bool
    kicking_team: int
    secs_remaining: int
    secondary_time: int
    teams: list          # two dicts: number, score, goalkeeper, penalties

def parse_control(buf: bytes) -> GameState:
    if len(buf) != CONTROL.size + 2 * TEAM.size:
        raise ValueError(f"length {len(buf)}")
    (hdr, version, packet, _players, _comp, stopped, _phase, state,
     set_play, first_half, kicking, secs, secondary) = CONTROL.unpack_from(buf)
    if hdr != b"RGme":
        raise ValueError("bad header")
    if version != 20:
        raise ValueError(f"version {version}")
    teams = []
    for k in range(2):
        f = TEAM.unpack_from(buf, CONTROL.size + k * TEAM.size)
        teams.append({"number": f[0], "goalkeeper": f[3], "score": f[4],
                      "budget": f[7], "penalty": list(f[8::3]),
                      "secs_left": list(f[9::3])})
    return GameState(packet, STATES[state], stopped == 1, set_play,
                     first_half == 1, kicking, secs, secondary, teams)

def encode_return(player, team, fallen, pose, ball_age, ball):
    """pose = (x mm, y mm, theta rad), ball = (x mm, y mm) robot-relative."""
    return RETURN.pack(b"RGrt", 4, player, team, int(fallen),
                       *pose, ball_age, *ball)
\end{verbatim}
\verb!f[8::3]! picks every third value starting at the first player's penalty byte: the $20$ players' penalties; \verb!f[9::3]! their remaining penalty seconds. The struct sizes are checked by Python: \texttt{CONTROL.size} is $18$, \texttt{TEAM.size} $70$, \texttt{RETURN.size} $32$.

I tested the parser with a packet built by \texttt{struct} (state \emph{playing}, $540$\,s remaining, our player $3$ penalised for $30$\,s, budget $1200$): all fields came out as built; a packet one byte short was rejected with \emph{length 157}, and one with a wrong header with \emph{bad header}. The return packet was $32$ bytes and starts with \texttt{RGrt}, version $4$, player $3$, team $12$. This test used packets built from the layout in the book, not packets from the real GameController; Step~3 closes that gap.

\section*{Into the blackboard}
A listener thread receives the broadcast on port $3838$ and keeps the newest valid \texttt{GameState} (with its arrival time) in a \texttt{LatestBuffer} (Chapter~\ref{k:ros}). The packet also tells the robot the address of the GameController (the sender of the broadcast); the status packet goes back there, unicast to port $3939$, once per second. In cognition, a module \texttt{GameControllerInput} provides \texttt{GameState} and \texttt{OurTeam} (which of the two team blocks is ours, by team number from \texttt{config/team.toml}, and our penalty).
\begin{verbatim}
def listen(buffer, port=3838):
    sock = socket.socket(socket.AF_INET, socket.SOCK_DGRAM)
    sock.setsockopt(socket.SOL_SOCKET, socket.SO_REUSEADDR, 1)
    sock.bind(("", port))
    while True:
        data, addr = sock.recvfrom(1024)
        try:
            buffer.put((parse_control(data), addr[0]))
        except ValueError as e:
            log.warning("bad GameController packet: %s", e)
\end{verbatim}

\section*{No GameController?}
In the lab there is often no GameController running. The robot must then behave sensibly: start in \emph{initial} (standing still), and accept a game state from robotcontrol instead (an extra \texttt{Request} type). When the GameController's packets stop during a game (Wi-Fi), the robot keeps the last state and does not invent a new one. Write down for your team what the current rules require here\guess.

\section*{Steps}
\begin{stepbox}{Step 1: the parser and tests}
\texttt{gamecontroller.py} as above, with the tests described.
\end{stepbox}
\begin{stepbox}{Step 2: listener, sender, module}
The thread, the status sender at $1$\,Hz (player number and team from the configuration; pose and ball from the blackboard; \emph{fallen} from Chapter~\ref{k:deploy}), and the \texttt{GameControllerInput} module.
\end{stepbox}
\begin{stepbox}{Step 3: against the real GameController}
Start the official GameController on the laptop with your team number, run the robot code with the fake robot on the same laptop, and click through \emph{ready}, \emph{set}, \emph{playing}, a penalty, a goal. Record the packets once with \verb!tcpdump -i any -w gc.pcap udp port 3838! and add a few of them as test files.
\end{stepbox}
\begin{stepbox}{Step 4: show it}
Add the game state, our penalty, the score and the time to robotcontrol's view.
\end{stepbox}

\begin{checkbox}
\begin{itemize}[nosep]
  \item robotcontrol shows every state change within one second of the click.
  \item The GameController's window shows your robot as connected (it receives the status packets).
  \item The parser tests pass with packets recorded from the real GameController.
\end{itemize}
\end{checkbox}

\begin{troublebox}
\begin{description}[leftmargin=0.6cm,nosep]
  \item[No packets] Broadcasts do not cross routers, and some Wi-Fi networks block them. Test on one laptop first; then check with \verb!tcpdump!.
  \item[\emph{version} errors] The GameController in use speaks another version than the one in the book. Update the layout from its header file.
  \item[Our team is never found] The team number in \texttt{config/team.toml} differs from the one chosen in the GameController.
\end{description}
\end{troublebox}

\begin{keyidea}[What to remember]
\begin{itemize}[nosep]
  \item \texttt{struct} formats express the packet layouts in one line; check length, header and version before trusting a packet.
  \item Listener in a thread, newest valid state in a buffer, a module puts it on the blackboard; status back once per second.
  \item Test with packets from the real GameController.
\end{itemize}
\end{keyidea}

\section*{Practice problems}
\begin{enumerate}[label=\textbf{\thechapter.\arabic*}]
  \item Compute the size of \verb!"<6B2H60B"! by hand. What would the size be without the \texttt{<} (native alignment), and why does that matter?
  \item Write \verb!our_team(state, team_number)! that returns the index ($0$ or $1$) of our team block, or raises an error.
  \item Player $3$ is penalised. Which index in \texttt{teams[k]["penalty"]} holds his or her penalty, and why?
  \item The status packet wants the ball in millimetres relative to the robot. Our \texttt{BallModel} has metres. Write the call to \verb!encode_return! for a ball at $(1.2, -0.3)$\,m seen $0.4$\,s ago and a pose of $(-2.5,1.0)$\,m, $0.3$\,rad.
  \item The packet number is one byte. After how many seconds does it wrap around at $2$ packets per second? How do you detect a lost packet across the wrap?
\end{enumerate}
